% !PRINTED none

\begin{center}
\gesperrt{Twenty-Second Chapter.}

\gesperrt{Abelian Fields and Cyclotomic Fields.}

\rule{3cm}{0.4pt}

\S.\,179.

\gesperrt{Simple Abelian Fields.}
\end{center}

We turn now to one of the most interesting applications of the theory of algebraic numbers.

\begingroup
\renewcommand{\thefootnote}{\arabic{footnote})}
It concerns the proof of the general theorem\footnote{\gesperrt{Kronecker} first stated this theorem, but published no complete proof of it. The first proof was made known by the author of the present work in Bd. 8 of the Acta Mathematica (1886). Most recently a second proof by \gesperrt{Hilbert} has been published, which rests upon the theorems developed in the nineteenth chapter (Nachrichten d. Gesellschaft d. Wissenschaften zu Göttingen, 1896, Heft 1).}:
\endgroup

I. \gesperrt{All number fields that are Abelian in the absolute domain of rationality are cyclotomic fields.}

Once we have proved this theorem, the investigations of the third chapter on cyclotomic fields gain a heightened interest, because it is thereby shown that on the path indicated there not only all cyclotomic fields, but all \gesperrt{Abelian} fields whatever, are found.

Let $\Omega = R(x)$ be an \gesperrt{Abelian} field of degree $m$, that is, a field which consists of all rational functions of a root $x$ of an irreducible \gesperrt{Abelian} equation of degree $m$, when the field $R$ of the rational numbers is taken as the domain of rationality. The \gesperrt{Galois} group $\mathfrak{G}$ of this equation, which is thus an \gesperrt{Abelian} group and has the same degree $m$ as the field $\Omega$, is also the group of the field $\Omega$.

If $x_1$ is a non-primitive number from $\Omega$, then $\Omega_1 = R(x_1)$ is likewise an \gesperrt{Abelian} field, which we call a proper divisor of $\Omega$, and whose degree is a proper divisor of the degree of $\Omega$. For if $\mathfrak{A}$ is the group to which the number $x_1$ belongs,
%CONT>>
